\chapter{Users and their needs}\label{ch:user-needs}
% Draft by Claude, edit freely. Everything in this chapter is a plan. Nothing has been done yet.
% Rewrite it in past tense with the results once the interviews and the email are done.

\fillin{This chapter is a plan. None of these studies have been done yet. Update it with what
we actually did and found.}

In the introduction we named two kinds of users: people who would build a station at home, and
the people who run the current pollen stations. Before we design more of the device, we want
to know what these two groups need from it. A cheap device is only useful if the data it gives
is something people want and trust.

We plan two small studies, one for each group.

\section{Interviews with people with pollen allergy}\label{sec:allergy-interviews}

\fillin{Planned, not done yet.}

\paragraph{Goal.} Find out how people with pollen allergy use pollen information today, and
whether a local count would change what they do.

\paragraph{Participants.} 3 to 4 people with pollen allergy.
\fillin{Who, and how we found them (for example friends, family or other students).}
This is a small group, so the results show possible needs, not what most people with allergy
think.

\paragraph{Method.} Semi-structured interviews of about 30 minutes, one person at a time. We
follow a list of questions (Appendix~\ref{app:interview-guide}) but ask follow-up questions when
something is interesting. With the participant's consent we record the audio, or take notes if
they prefer. We look for themes that come up in several interviews, using thematic
analysis~\cite{Braun2006}.

\paragraph{Topics.}
\begin{itemize}
  \item Where they get pollen information today, and how often they check it.
  \item What they do with it (medicine, staying inside, planning exercise or trips).
  \item Whether the official count matches how they feel where they live.
  \item How local and how frequent the information would need to be to help them.
  \item Whether they would trust a count from a cheap station run by a private person.
  \item Whether they would build or host a station themselves, and what would stop them.
\end{itemize}

\paragraph{Consent and data.} Allergy symptoms are health information. We ask for written
consent, keep the recordings only until the notes are written, and do not use names in the
report. \fillin{Check with the supervisor if AAU needs anything more for this.}

\section{Written questions to Astma-Allergi Danmark}\label{sec:expert-email}

\fillin{Planned, not done yet.}

\paragraph{Goal.} Learn how the people who count pollen today work, and what they would need
before they could use data from cheap stations.

\paragraph{Method.} A meeting is hard to arrange, so we send a short email with written
questions (Appendix~\ref{app:expert-questions}). This is a written interview: the answers are
shorter than in a conversation and we can only ask follow-up questions in a new email, but it
takes little of their time, which makes an answer more likely.
\fillin{Who we sent it to, when, and whether they answered.}

\paragraph{Topics.}
\begin{itemize}
  \item How the daily count is made, and how much work it takes.
  \item What accuracy and time resolution a count needs to be useful for a forecast.
  \item What they would need to see before they trust an automated count.
  \item Whether a network of cheap stations would be useful to them, and who would look after
  the stations.
\end{itemize}

\section{What we learned}\label{sec:user-needs-results}

\fillin{Results from the interviews and the email, and what we change in the design because of
them. Write this when the studies are done.}
